\chapterbackmatter{Weinberg Exercises}

\begin{enumerate}
\WeinbergExercise{1} Suppose that observer $\mathcal{O}$ sees a $W$-boson (spin one and mass
  $m\neq0$) with momentum $\mathbf{p}$ in the $y$-direction and spin
  $z$-component $\sigma$. A second observer $\mathcal{O}'$ moves relative to
  the first with velocity $\mathbf{v}$ in the $z$-direction. How does
  $\mathcal{O}'$ describe the $W$ state?

\WeinbergExercise{2} Suppose that observer $\mathcal{O}$ sees a photon with momentum
  $\mathbf{p}$ in the $y$-direction and polarization vector in the
  $z$-direction. A second observer $\mathcal{O}'$ moves relative to the first
  with velocity $\mathbf{v}$ in the $z$-direction. How does $\mathcal{O}'$
  describe the same photon?

\WeinbergExercise{3} Derive the commutation relations for the generators of the Galilean
  group directly from the group multiplication law (without using our
  results for the Lorentz group). Include the most general set of central
  charges that cannot be eliminated by redefinition of the group generators.

\WeinbergExercise{4} Show that the operators $P_\mu P^\mu$ and $W_\mu W^\mu$ commute with all
  Lorentz transformation operators $U(\Lambda,a)$, where
  $W_\mu\equiv\epsilon_{\mu\nu\rho\lambda}J^{\nu\rho}P^\lambda$.

\WeinbergExercise{5} Consider physics in two space and one time dimensions, assuming
  invariance under a `Lorentz' group $SO(2,1)$. How would you describe the
  spin states of a single massive particle? How do they behave under Lorentz
  transformations? What about the inversions P and T?

\WeinbergExercise{6} As in Problem 5, consider physics in two space and one time dimensions,
  assuming invariance under a `Lorentz' group $O(2,1)$. How would you describe
  the spin states of a single massless particle? How do they behave under
  Lorentz transformations? What about the inversions P and T?
\end{enumerate}
